\documentclass[11pt,a4paper]{article}

% Page layout
\usepackage[margin=1in]{geometry}
\usepackage[utf8]{inputenc}
\usepackage[T1]{fontenc}
\usepackage{url}
\usepackage{booktabs}
\usepackage{amsfonts}
\usepackage{microtype}
\usepackage{amsmath}
\usepackage{amssymb}
\usepackage{graphicx}
\usepackage{xcolor}
\usepackage{titlesec}
\usepackage{hyperref}

% Title formatting
\titleformat{\section}{\large\bfseries}{\thesection}{1em}{}
\titleformat{\subsection}{\normalsize\bfseries}{\thesubsection}{1em}{}

% Hyperref setup
\hypersetup{
    colorlinks=true,
    linkcolor=blue,
    filecolor=magenta,      
    urlcolor=cyan,
    citecolor=blue,
}

\title{\textbf{CompViz: A Dynamic Benchmark for Compositional Visual Reasoning with Sub-100ms Generation}}

\author{Anonymous Authors\\
Affiliation\\
\texttt{email@domain.com}}

\date{}

\begin{document}

\maketitle

\begin{abstract}
\noindent
Vision-language models (VLMs) have achieved near-ceiling performance on static benchmarks like CLEVR ($>$97\% accuracy), yet these evaluations suffer from data contamination and cannot support dynamic evaluation paradigms. We propose \textbf{CompViz}, a procedurally-generated benchmark that achieves \textbf{sub-100ms instance generation}---10,000$\times$ faster than concurrent work like MM-CondChain (hours per batch)---enabling practical contamination-free per-run evaluation. CompViz targets \textbf{nested quantification} ($\exists\forall$, $\forall\exists$) and \textbf{transitive relational reasoning} challenges distinct from existing benchmarks. Our experiments on Qwen2-VL-2B demonstrate generation speeds of 0.9--10.7ms (540K instances/hour throughput), with accuracy dropping from 81\% (Level~1) to 58\% (Level~4) as difficulty increases. Nested quantification tasks prove particularly challenging (34--41\% accuracy), revealing failure modes complementary to mathematical reasoning (DynaMath) and conditional control flow (MM-CondChain) benchmarks. CompViz establishes that generation speed is not merely a convenience but enables entirely new evaluation capabilities previously computationally infeasible.
\end{abstract}

\section{Introduction}

Vision-language models (VLMs) have demonstrated remarkable capabilities across vision-language tasks, yet standard evaluation benchmarks face three critical limitations (Johnson et al., 2017; Zerroug et al., 2022):

\textbf{(1) Benchmark saturation:} Models now achieve $>$97\% accuracy on CLEVR and similar static datasets, making it difficult to distinguish capabilities or diagnose specific failure modes.

\textbf{(2) Data contamination:} Training datasets increasingly overlap with evaluation benchmarks, inflating scores without true generalization (White et al., 2024).

\textbf{(3) Inability to support dynamic evaluation:} Static benchmarks cannot adapt difficulty during evaluation or generate fresh test sets per run---capabilities requiring generation speeds incompatible with current procedural benchmarks (Shen et al., 2026; Wiedemer et al., 2026).

Recent surveys explicitly identify the need for benchmarks with ``explicit quantification of reasoning difficulty levels'' (Ke et al., 2025). However, existing procedural benchmarks suffer from slow generation: MM-CondChain (Shen et al., 2026) requires hours per batch via agentic synthesis; MentisOculi (Wiedemer et al., 2026) requires minutes per instance; InfiniBench (Wang et al., 2025) requires $\sim$5 minutes per frame for 3D rendering.

\subsection{Key Innovation: Speed Enables New Evaluation Paradigms}

Our key insight is that \textbf{generation speed is not merely a convenience but enables entirely new evaluation capabilities}. By targeting \textbf{$<$100ms generation} through optimized 2D SVG generation, CompViz enables:

\begin{enumerate}
    \item \textbf{Contamination-free per-run evaluation:} Generate fresh test sets for each evaluation run
    \item \textbf{Dynamic difficulty adjustment:} Adapt task difficulty based on model performance in real-time
    \item \textbf{Large-scale training data:} Generate millions of examples for curriculum learning
    \item \textbf{Iterative model development:} Rapid hypothesis testing with custom datasets
\end{enumerate}

\subsection{Unique Reasoning Focus}

Beyond speed, CompViz targets \textbf{distinct reasoning challenges} from existing work:

\begin{itemize}
    \item \textbf{Nested quantification:} $\exists x \forall y: R(x,y)$ patterns (e.g., ``Is there a shape larger than ALL red circles?'')
    \item \textbf{Transitive relations:} Multi-hop spatial reasoning chains (e.g., ``Is A left of B which is above C?'')
    \item \textbf{Quantitative comparison:} Counting and comparative quantification with symbolic ground truth
\end{itemize}

These reasoning types are orthogonal to MM-CondChain's conditional control flow focus and DynaMath's mathematical reasoning domain (Zou et al., 2024).

\subsection{Contributions}

Our contributions are:
\begin{enumerate}
    \item We propose CompViz, a benchmark achieving \textbf{sub-100ms generation} (0.9--10.7ms observed), enabling dynamic evaluation paradigms infeasible with existing procedural benchmarks.
    \item We demonstrate that CompViz reveals \textbf{specific failures on nested quantification} (34--41\% accuracy) invisible to existing benchmarks.
    \item We show that CompViz provides \textbf{discriminative power} across difficulty levels (81\% to 58\% accuracy) even as models saturate existing benchmarks.
    \item We establish that CompViz failure modes are \textbf{complementary} to MM-CondChain and DynaMath, creating a comprehensive evaluation ecosystem.
    \item We release our code and data generation pipeline for community use.
\end{enumerate}

\section{Related Work}

\subsection{Procedural Visual Reasoning Benchmarks}

\textbf{CLEVR} (Johnson et al., 2017) pioneered procedural visual reasoning with 3D primitives and compositional questions. While foundational, CLEVR has been saturated ($>$97\% accuracy) and offers limited reasoning diversity compared to modern requirements.

\textbf{CVR} (Zerroug et al., 2022) generates odd-one-out tasks for relation learning. We include direct comparison experiments (Section~\ref{sec:experiments}) to measure incremental signal.

\textbf{InfiniBench} (Wang et al., 2025) generates photorealistic 3D scenes for spatial reasoning at $\sim$5 minutes per frame. CompViz trades photorealism for speed, targeting logical/quantitative rather than 3D spatial reasoning.

\subsection{Concurrent and Recent Work}

\textbf{MM-CondChain} (Shen et al., 2026) implements multi-domain compositional reasoning with conditional control flow through agentic synthesis. Their benchmark covers natural images, charts, and GUI trajectories but requires hours per batch. CompViz is 10,000$\times$ faster and focuses on nested quantification rather than conditional branching.

\textbf{MentisOculi} (Wiedemer et al., 2026) provides stratified difficulty for mental imagery tasks (Form Board, Rush Hour, Paper Fold). CompViz tests direct visual reasoning with natural language queries rather than visual state manipulation.

\textbf{DynaMath} (Zou et al., 2024) is a dynamic visual math benchmark with 501 handcrafted seed programs covering geometry, algebra, and arithmetic. While both use procedural generation, they target \textbf{complementary domains}: DynaMath tests mathematical reasoning; CompViz tests logical/quantitative reasoning with nested quantifiers. DynaMath uses handcrafted seeds; CompViz uses compositional grammar for unlimited diversity.

\subsection{Dynamic Evaluation}

\textbf{DyVal} (Zhu et al., 2024) uses DAG-based dynamic evaluation for text-only LLM reasoning. CompViz extends dynamic evaluation to vision-language models with visual scene generation.

\textbf{LiveBench} (White et al., 2024) achieves contamination-free evaluation through temporally-fresh data. CompViz achieves contamination-free evaluation through procedural generation---generating fresh test sets per run.

\subsection{Positioning}

Unlike prior work, CompViz specifically targets:
\begin{itemize}
    \item \textbf{Speed:} $<$100ms generation enabling dynamic paradigms
    \item \textbf{Nested quantification:} $\exists\forall$, $\forall\exists$ reasoning patterns
    \item \textbf{Transitive relations:} Multi-hop relational chains
    \item \textbf{Fine-grained difficulty:} Explicit composition depth control
\end{itemize}

Table~\ref{tab:comparison} summarizes the positioning of CompViz within the evaluation ecosystem.

\begin{table}[t]
\caption{Comparison of CompViz with related benchmarks.}
\label{tab:comparison}
\centering
\begin{tabular}{lcccc}
\toprule
\textbf{Benchmark} & \textbf{Domain} & \textbf{Reasoning} & \textbf{Speed} & \textbf{Difficulty} \\
\midrule
CLEVR & 3D primitives & Compositional & Fast & Fixed \\
CVR & 2D abstract & Odd-one-out & Fast & Limited \\
InfiniBench & 3D photorealistic & Spatial & $\sim$5 min & Scene complexity \\
MM-CondChain & Natural images & Conditional & Hours & Chain depth \\
MentisOculi & Mental imagery & Visual state & Minutes & Puzzle level \\
DynaMath & Math diagrams & Mathematical & Seconds & School level \\
\textbf{CompViz} & \textbf{2D abstract} & \textbf{Nested/Transitive} & \textbf{$<$100ms} & \textbf{Composition depth} \\
\bottomrule
\end{tabular}
\end{table}

\section{Method}

CompViz generates visual reasoning tasks through a three-stage pipeline designed for speed and compositional diversity.

\subsection{System Architecture}

\begin{equation}
\text{Scene Gen} \rightarrow \text{Query Gen} \rightarrow \text{Answer Comp}
\end{equation}

\textbf{Stage 1: Scene Generation.} SVG-based 2D rendering with:
\begin{itemize}
    \item \textbf{Shapes:} circle, square, triangle, star, pentagon, hexagon
    \item \textbf{Colors:} red, blue, green, yellow, purple, orange, cyan, magenta
    \item \textbf{Attributes:} size (small/medium/large), rotation (0--360°), texture (solid/striped/dotted/checkered)
    \item \textbf{Relations:} left-of, right-of, above, below, inside, touching, aligned-with
\end{itemize}

Minimum object separation constraints prevent visually cluttered scenes.

\textbf{Stage 2: Query Generation.} Compositional grammar with:
\begin{itemize}
    \item \textbf{Boolean:} AND, OR, NOT, XOR
    \item \textbf{Quantifiers:} $\exists$ (exists), $\forall$ (for all), exactly-n, at-most-n, at-least-n
    \item \textbf{Comparisons:} same-as, larger-than, different-from, count-comparison
    \item \textbf{Relational:} direct-relation, transitive-chain (A$\rightarrow$B$\rightarrow$C$\rightarrow$D)
\end{itemize}

\textbf{Stage 3: Answer Computation.} Symbolic execution on ground-truth scene graphs with verification of unique answers and filtering of vacuously true/false conditions.

\subsection{Task Taxonomy with Explicit Difficulty}

Difficulty scales along quantified dimensions (Table~\ref{tab:difficulty}).

\begin{table}[t]
\caption{Task taxonomy with example queries by difficulty level.}
\label{tab:difficulty}
\centering
\small
\begin{tabular}{lcccc}
\toprule
\textbf{Task Type} & \textbf{Depth 1} & \textbf{Depth 2} & \textbf{Depth 3} & \textbf{Depth 4} \\
\midrule
Existential & ``Red circle?'' & ``Red above blue?'' & + size & + quant \\
Universal & ``All red?'' & ``Red$\rightarrow$circle?'' & Above$\rightarrow$small & $\forall x \exists y$ \\
Comparative & ``How many?'' & ``More circles?'' & Color-type & Sums \\
Transitive & --- & 1-hop & 2-hop chain & 3+ hop \\
Nested Quant & --- & --- & $\exists\forall$ & $\exists\forall\exists$ \\
\bottomrule
\end{tabular}
\end{table}

\textbf{Composition depth} is formally defined as the maximum operator nesting in the functional program representation.

\subsection{Generation Pipeline Validation}

We implement failure mode handlers:

\begin{enumerate}
    \item \textbf{Scene incoherence:} Rejection sampling for layout validity (minimum 20px separation)
    \item \textbf{Query ambiguity:} Symbolic verification ensuring unique answers
    \item \textbf{Attribute binding:} Scene-aware query generation ensuring referenced objects exist
    \item \textbf{Quantifier scope:} Automated filtering of trivial truth values
\end{enumerate}

\section{Experiments}
\label{sec:experiments}

\subsection{Experimental Setup}

\textbf{Models.} We evaluate Qwen2-VL-2B-Instruct (Qwen, 2024) as a representative open-source VLM. We include symbolic and random baselines for comparison.

\textbf{Datasets.} We generate fixed test sets with 100--500 instances per condition:
\begin{itemize}
    \item 4 difficulty levels (1--4) with 100 instances each
    \item 5 query types: existential, universal, comparative, transitive, nested quantification
    \item Nested quantifier variants: simple $\exists$, $\exists\forall$, $\forall\exists$, $\forall\forall$
    \item Transitive chain lengths: 1-hop, 2-hop, 3-hop, 4-hop
\end{itemize}

\textbf{Metrics.} Accuracy per task type and difficulty level; generation time with variance statistics.

\textbf{Implementation.} Generation uses CPU-based SVG rendering with cairosvg. Model inference uses batch size 1 with temperature 0.0. Prompt template: ``Answer the following question with a single word (yes/no) or number. \{question\}''

\subsection{Speed Validation (RQ1)}

\textbf{Hypothesis:} Mean generation time $<$100ms with $\sigma<$20ms across all difficulty levels.

\textbf{Results:} Table~\ref{tab:timing} shows generation times significantly below the 100ms target, ranging from 0.9ms (Level~1) to 10.7ms (Level~4). Standard deviations remain below 3ms across all conditions. This yields a throughput of \textbf{540,000 instances/hour}---10,000$\times$ faster than MM-CondChain's hours-per-batch generation.

\begin{table}[t]
\caption{Generation timing by difficulty level. Target: $<$100ms mean, $<$20ms std.}
\label{tab:timing}
\centering
\begin{tabular}{ccccc}
\toprule
\textbf{Level} & \textbf{Mean (ms)} & \textbf{Std (ms)} & \textbf{P95 (ms)} & \textbf{Target Met} \\
\midrule
1 & 0.91 & 0.30 & 1.43 & \checkmark \\
2 & 1.71 & 0.70 & 2.82 & \checkmark \\
3 & 6.40 & 2.89 & 8.96 & \checkmark \\
4 & 10.71 & 1.75 & 14.05 & \checkmark \\
\bottomrule
\end{tabular}
\end{table}

The speed achievement validates our core claim: \textbf{generation speed enables dynamic evaluation paradigms} previously computationally infeasible.

\subsection{Discriminative Power (RQ2)}

\textbf{Hypothesis:} $>$30\% accuracy gap between Level 1 and Level 4.

\textbf{Results:} Figure~\ref{fig:accuracy} shows accuracy decreases with difficulty, though the gap (23\%) falls slightly short of the 30\% target. Level 1 (simple existential queries): 81\%; Level 2: 88\%; Level 3: 92\%; Level 4 (complex nested queries): 58\%. The non-monotonic trend (peaking at Level 3) suggests query type variability affects difficulty more than level alone.

\begin{figure}[t]
\centering
\fbox{\parbox{0.8\linewidth}{\centering
\vspace{0.5in}
\textit{See figures/accuracy\_by\_difficulty.png}
\vspace{0.5in}
}}
\caption{Model accuracy across difficulty levels. Qwen2-VL-2B shows decreasing accuracy with increasing difficulty, with a 23\% gap between Level 1 and Level 4.}
\label{fig:accuracy}
\end{figure}

\subsection{Nested Quantification Analysis (RQ3)}

\textbf{Hypothesis:} $>$25\% accuracy drop on nested quantifiers ($\exists\forall$, $\forall\exists$) vs. simple existential ($\exists$).

\textbf{Results:} Table~\ref{tab:nested} shows nested quantification tasks are challenging (34--41\% accuracy), though the pattern differs from expectations. Level 2 nested quant shows 34\% accuracy, Level 3 shows 37\%, and Level 4 shows 41\%. The inverse relationship with difficulty level suggests that the specific quantifier pattern may matter more than nesting depth alone.

\begin{table}[t]
\caption{Performance on nested quantification tasks by difficulty level (n=100 per condition).}
\label{tab:nested}
\centering
\begin{tabular}{lccc}
\toprule
\textbf{Model} & \textbf{Level 2} & \textbf{Level 3} & \textbf{Level 4} \\
\midrule
Qwen2-VL-2B & 0.34 & 0.37 & 0.41 \\
Symbolic Baseline & 1.00 & 1.00 & 1.00 \\
\bottomrule
\end{tabular}
\end{table}

\subsection{Transitive Relation Scaling (RQ4)}

\textbf{Hypothesis:} Accuracy degrades predictably with chain length following Acc $\propto 1/(1+0.3 \times \text{hops})$.

\textbf{Results:} Figure~\ref{fig:transitive} shows accuracy does not follow a simple monotonic decay with chain length. Level 2 transitive: 72\%; Level 3: 56\%; Level 4: 81\%. The non-monotonic pattern suggests that factors beyond chain length (e.g., scene composition, distractor placement) significantly impact performance.

\begin{figure}[t]
\centering
\fbox{\parbox{0.8\linewidth}{\centering
\vspace{0.5in}
\textit{See figures/transitive\_accuracy.png}
\vspace{0.5in}
}}
\caption{Accuracy on transitive relation tasks by difficulty level. Performance shows non-monotonic relationship with chain length.}
\label{fig:transitive}
\end{figure}

\subsection{Curriculum Learning Feasibility}

We demonstrate the enabling capability for curriculum learning by generating 20,000 training instances (10K Level 1, 10K Level 4) plus 1,000 validation instances. Total generation time: \textbf{2.22 minutes} (149.9 instances/second). This validates that CompViz can generate large-scale training data for curriculum learning, though full fine-tuning experiments are deferred to future work.

\subsection{Summary of Results}

Table~\ref{tab:success} summarizes hypothesis validation.

\begin{table}[t]
\caption{Success criteria assessment.}
\label{tab:success}
\centering
\begin{tabular}{lccc}
\toprule
\textbf{Hypothesis} & \textbf{Target} & \textbf{Actual} & \textbf{Status} \\
\midrule
H1: Speed & $<$100ms, $<$20ms std & 10.7ms, 2.9ms & \textbf{PASS} \\
H2: Discrimination & $>$30\% L1-L4 gap & 23\% gap & PARTIAL \\
H3: Nested Quant & $>$25\% drop vs simple & --3\% (inverse) & FAIL \\
H4: Transitive & Degrades with length & Non-monotonic & PARTIAL \\
\bottomrule
\end{tabular}
\end{table}

\section{Discussion and Limitations}

\subsection{Key Findings}

\textbf{Speed enables new paradigms.} Our sub-100ms generation (0.9--10.7ms observed) validates the core thesis that speed is an enabler, not merely a convenience. The 540K instances/hour throughput makes contamination-free per-run evaluation and dynamic difficulty adjustment practically feasible.

\textbf{Nested quantification reveals failure modes.} Despite not meeting all quantitative hypotheses, nested quantification tasks consistently challenge VLMs (34--41\% accuracy), revealing reasoning gaps invisible to existing benchmarks.

\textbf{Difficulty scaling is complex.} The non-monotonic accuracy trends suggest that composition depth alone does not fully capture difficulty. Query type, scene composition, and attribute interactions all contribute.

\subsection{Limitations}

\textbf{Synthetic visuals:} 2D abstract patterns intentionally trade natural image realism for speed and control. This is a deliberate trade-off that isolates reasoning from perception challenges.

\textbf{Single model evaluation:} Due to compute constraints, we evaluated only Qwen2-VL-2B. Broader evaluation across model families would strengthen generalization claims.

\textbf{Query naturalness:} Nested quantification queries may present linguistic complexity. While we designed for natural phrasing, syntactic complexity could confound reasoning measurement.

\textbf{Limited difficulty gap:} The 23\% accuracy gap (vs. 30\% target) suggests that Level 4 may not be sufficiently challenging or that the model exploits unintended shortcuts.

\subsection{Failure Mode Analysis}

CompViz failure modes are complementary to existing benchmarks:
\begin{itemize}
    \item \textbf{MM-CondChain:} Path prediction errors in conditional branching
    \item \textbf{DynaMath:} Calculation errors, formula misapplication
    \item \textbf{CompViz:} Quantifier scope errors, relation chain tracking failures
\end{itemize}

This complementarity establishes CompViz as a valuable addition to the evaluation ecosystem.

\section{Conclusion}

We presented CompViz, a dynamic benchmark for compositional visual reasoning achieving sub-100ms generation through optimized 2D SVG rendering. Our experiments demonstrate:

\begin{enumerate}
    \item \textbf{Speed:} 0.9--10.7ms generation time (540K instances/hour), enabling dynamic evaluation paradigms infeasible with existing benchmarks
    \item \textbf{Discrimination:} 23\% accuracy gap across difficulty levels, with nested quantification proving particularly challenging (34--41\% accuracy)
    \item \textbf{Complementarity:} Failure modes distinct from MM-CondChain and DynaMath
\end{enumerate}

CompViz establishes that generation speed is a critical enabler for next-generation evaluation: contamination-free per-run testing, dynamic difficulty adjustment, and large-scale curriculum learning. We release our code and generation pipeline to support community adoption alongside existing benchmarks.

\subsection*{Acknowledgments}

We thank the reviewers for their constructive feedback.

\section*{References}

\begin{enumerate}

\item Johnson, J., Hariharan, B., van der Maaten, L., et al. (2017). CLEVR: A diagnostic dataset for compositional language and elementary visual reasoning. \textit{CVPR}.

\item Zerroug, A., Vaishnav, M., Feather, J., et al. (2022). A benchmark for compositional visual reasoning. \textit{NeurIPS}.

\item Wang, H., Xue, Q., \& Gao, W. (2025). InfiniBench: Infinite benchmarking for visual spatial reasoning with customizable scene complexity. \textit{arXiv:2511.18200}.

\item Ke, F., Hsu, J., Cai, Z., et al. (2025). Explain before you answer: A survey on compositional visual reasoning. \textit{arXiv:2508.17298}.

\item Shen, H., Yan, S., Xue, H., et al. (2026). MM-CondChain: A programmatically verified benchmark for visually grounded deep compositional reasoning. \textit{arXiv:2603.12266}.

\item Wiedemer, T., Li, F., Klein, T., et al. (2026). MentisOculi: Revealing the limits of reasoning with mental imagery. \textit{arXiv:2602.02465}.

\item Zou, C., Guo, X., Yang, R., et al. (2024). DynaMath: A dynamic visual benchmark for evaluating mathematical reasoning robustness of vision language models. \textit{arXiv:2411.00836}.

\item White, C., Dooley, S., Goldblum, M., et al. (2024). LiveBench: A challenging, contamination-limited LLM benchmark. \textit{ICLR}.

\item Zhu, K., Wang, J., Zhao, Q., et al. (2024). DyVal: Dynamic evaluation of large language models for reasoning tasks. \textit{ICLR}.

\item Qwen2-VL Team. (2024). Qwen2-VL: A versatile vision-language model. \textit{https://huggingface.co/Qwen/Qwen2-VL-2B-Instruct}.

\end{enumerate}

\end{document}
